\documentclass[11pt]{article}

\usepackage[margin=1in]{geometry}
\usepackage{amsmath,amssymb}
\usepackage{graphicx}
\usepackage{booktabs}
\usepackage{xcolor}
\usepackage{algorithm}
\usepackage{algpseudocode}
\usepackage[colorlinks=true,linkcolor=blue!50!black,citecolor=blue!50!black,urlcolor=blue!50!black]{hyperref}

\graphicspath{{figures/}}

\newcommand{\Mmin}{M_{\min}}
\newcommand{\Mmax}{M_{\max}}
\newcommand{\ML}{M_{\mathrm{L}}}
\newcommand{\dM}{\Delta M_{\mathrm{B}}}
\newcommand{\Dres}{D}
\newcommand{\Dref}{D_{\mathrm{ref}}}
\newcommand{\Ltot}{L}
\newcommand{\fas}{f_{\mathrm{as}}}
\newcommand{\EL}{E_{\mathrm{L}}}
\newcommand{\EO}{E_{\mathrm{O}}}
\newcommand{\nL}{n_{\mathrm{L}}}
\newcommand{\nO}{n_{\mathrm{O}}}
\newcommand{\lamL}{\lambda_{\mathrm{L}}}
\newcommand{\lamO}{\lambda_{\mathrm{O}}}

\title{Loading-limited large earthquakes with triggered small-event sequences:\\
a moment-budget/Omori split for the KES stochastic earthquake simulator}
\author{Brendan J. Meade}
\date{\today}

\begin{document}
\maketitle

\begin{abstract}
\noindent
KES generates synthetic earthquake catalogs on a discretized strike-slip fault from
statistical-mechanics ingredients: tectonic loading of a slip-deficit field, an event
rate set by moment balance with a Dieterich-style stress-shadow memory, maximum-entropy
(MaxEnt) hypocenter selection, and magnitude--area rupture scaling with slip clipped to
the locally available deficit. In its classical configuration every event, large or
small, is an independent draw from a truncated Gutenberg--Richter (G-R) law, and moment
balance then fixes the fault-wide $M\ge 5$ rate near $1.5\,\mathrm{yr^{-1}}$: the
interseismic period is as active as any other, unlike natural catalogs, in which small
events concentrate into aftershock sequences. This note documents a variant that splits
the event population into two physically distinct classes tied together by a single
moment budget. \emph{Loading events} are nucleated by the accumulated moment deficit and
draw magnitudes only from the top of the distribution, $M \in [\ML,\Mmax] = [6, 7.5]$;
\emph{Omori children} are triggered by previous events, occur at Omori--Utsu rates
proportional to a productivity scaled tenfold above the classical setting, and draw
magnitudes from the full G-R law capped $\dM = 1.2$ magnitude units below their parent
(B\aa th's law). Moment balance closes over the triggering cascade, which yields the
loading rate in closed form and leaves the long-run seismic coupling at unity without
retuning. In five $3000$-yr realizations the model produces $1.55$ $M\ge5$ events per
year of which $89\%$ are sequence members, a bursty catalog (interevent coefficient of
variation $2.4$; one-year Fano factor $8.6$), quiet interseismic periods at small
magnitudes, and a post-mainshock to pre-mainshock count ratio of $\sim3.6$ in the decade
around the largest events --- the ``nature-like'' texture that the classical
configuration lacks.
\end{abstract}

\section{Motivation}

Moment balance is a powerful constraint on stochastic seismicity models, but combined
with a single G-R magnitude law it has an under-appreciated corollary. If every event is
an independent draw from a truncated G-R density $p(M)\propto 10^{-bM}$ on
$[\Mmin,\Mmax]$, the mean geometric moment per event
$\bar m = \mathbb{E}[m(M)]$ is fixed, and balancing the aseismically corrected loading
rate $(1-\fas)\Ltot$ requires an event rate
\begin{equation}
  \bar\lambda \;=\; \frac{(1-\fas)\,\Ltot}{\bar m}
  \;\approx\; \frac{0.8 \times 7.67\times10^{7}\ \mathrm{m^3\,yr^{-1}}}
  {3.98\times10^{7}\ \mathrm{m^3}} \;\approx\; 1.5\ \mathrm{yr^{-1}}
\label{eq:pinned}
\end{equation}
for the reference fault used here ($b=1$, $\Mmin=5$, $\Mmax=7.5$). No knob that merely
redistributes events in time --- deeper stress shadows, deficit-dependent rate factors,
or added aftershock productivity --- can lower this long-run background: the shadow is
compensated by recalibration, and aftershocks are so moment-cheap that they add events
almost without displacing any. Natural catalogs, and quasi-static depinning models run
as a comparison target for KES, instead show \emph{quiet} interseismic periods: small
events cluster into aftershock sequences attached to the large ones.

The variant described here achieves that texture while keeping every conservation
property of the original model. The magnitude law is split by \emph{origin}: the
deficit-driven (loading) process nucleates only events large enough to matter for the
moment budget, $M \ge \ML = 6$, and all smaller events are generated by Omori--Utsu
triggering off previous events. The moment budget is then closed over the triggering
cascade (Section~\ref{sec:budget}), which determines the loading rate
($0.17\,\mathrm{yr^{-1}}$, one $M\ge6$ per $\sim6$ yr) and, for a productivity
multiplier $\kappa_K = 10$, restores the classical total rate of
$\sim1.5\ M{\ge}5\,\mathrm{yr^{-1}}$ --- now with nine of every ten events belonging to
a sequence (Figure~\ref{fig:lollipop}).

\begin{figure}[t]
  \centering
  \includegraphics[width=\textwidth]{lollipop_run00.png}
  \caption{One thousand years of the model (run 00, seed 42; the 1000-yr window that
  places the run's largest event at 35\%). Circles are loading-origin events ($M\ge\ML=6$); red-edged
  squares are Omori children (all $M<6.3$ here, B\aa th-capped $1.2$ units below their
  parents). Stem base at $M=5$; marker area $\propto 3.2^{\,M-5}$. Small-magnitude
  activity occurs almost exclusively in bursts attached to the larger events; the
  interseismic intervals are quiet at $M\approx5$.}
  \label{fig:lollipop}
\end{figure}

\section{Model definition}
\label{sec:model}

\subsection{Fault, loading, and the slip-deficit field}

The fault is a planar vertical strike-slip surface, $200\ \mathrm{km}$ along strike
$\times\ 25\ \mathrm{km}$ down dip, discretized into $N = 400\times50 = 20{,}000$
square elements of side $0.5$~km and area $a = 2.5\times10^{5}\ \mathrm{m^2}$. Element
$i$ accumulates slip deficit $s_i(t)$ (meters) at rate $v_i$: a uniform background of
$10\ \mathrm{mm\,yr^{-1}}$ plus one Gaussian loading pulse ($20\ \mathrm{mm\,yr^{-1}}$
amplitude, $\sigma = 25$~km, centered at $x=100$~km, $z=0$). All bookkeeping uses
\emph{geometric moment} $m_i = a\,s_i$ ($\mathrm{m^3}$); seismic moment is
$M_0 = \mu m$ with $\mu = 30$~GPa, and moment magnitude follows
$M_0 = 10^{1.5M + 9.05}$~N\,m. The total loading rate is
$\Ltot = a\sum_i v_i = 7.67\times10^{7}\ \mathrm{m^3\,yr^{-1}}$, and the fault-wide
reservoir is $\Dres(t) = a\sum_i s_i(t)$. The reservoir is initialized (and regulated,
below) at $\Dref = T_{\mathrm{ref}}\Ltot$ with $T_{\mathrm{ref}} = 250$~yr of stored
loading, deep enough that no allowed rupture is capacity-limited (an $M7.5$ requires
$134$~yr of storage).

\subsection{Event rate: memory stress shadow with reservoir feedback}
\label{sec:rate}

The loading-origin event rate is
\begin{equation}
  \lamL(t) \;=\; \lambda_0\, c(t)\, \langle h(t)\rangle
  \left(\frac{\Dres(t)}{\Dref}\right)^{\nu},
  \qquad \nu = \tfrac12 ,
\label{eq:rate}
\end{equation}
where $\langle h\rangle = N^{-1}\sum_i h_i$ is the mean of a per-element Dieterich-style
stress-shadow state. Each element carries $\xi_i \ge 1$ with $h_i = 1/\xi_i$; a rupture
placing coseismic slip $u_i$ on element $i$ applies a stress step
$\ln\xi_i \mathrel{+}= f\,u_i/(v_i t_a)$, and between events
$\xi \to 1 + (\xi-1)e^{-\delta t/t_a}$, with $f = 0.5$ and $t_a = 10$~yr. The
construction obeys the exact identity $\int (1-h_i)\,dt = f u_i / v_i$: a rupture
shadows its own patch for a fraction $f$ of the time needed to reload the slip it
released, independent of $t_a$ and additive across ruptures. The long-run mean shadow
under moment balance is $\langle h\rangle_{\mathrm{ref}} = 1 - f(1-\fas) = 0.6$, which
fixes $\lambda_0 = \bar\lambda_{\mathrm{L}}/\langle h\rangle_{\mathrm{ref}}$ once the
balanced rate $\bar\lambda_{\mathrm{L}}$ is known (Section~\ref{sec:budget}). The factor
$(\Dres/\Dref)^{\nu}$ is the proportional term and
$\dot c = g\,(\Dres/\Dref - 1)$, $g = 2\times10^{-4}\ \mathrm{yr^{-1}}$,
$c\in[0.5,2]$, the integral term of a proportional--integral controller that pins the
reservoir to $\Dref$ (damping ratio $\zeta = \nu/2\sqrt{g\,T_{\mathrm{ref}}} = 1.12$);
in practice $c$ stays within $\pm17\%$ of unity.

\subsection{The origin split and the two magnitude laws}
\label{sec:split}

Every event belongs to one of two classes.

\paragraph{Loading events} occur at rate $\lamL(t)$ and draw magnitude from the
truncated G-R density restricted to the top of the distribution,
\begin{equation}
  p_{\mathrm{L}}(M) \;\propto\; 10^{-bM}, \qquad M \in [\ML,\ \Mmax] = [6,\ 7.5],
  \quad b = 1 .
\label{eq:loadlaw}
\end{equation}
Their expected geometric moment is
$\EL = \mathbb{E}_{p_{\mathrm{L}}}[m] = 3.57\times10^{8}\ \mathrm{m^3}$, a factor $9.0$
above the full-range G-R mean --- this single change is what depresses the loading rate
from $1.5$ to $0.17\ \mathrm{yr^{-1}}$ by Eq.~(\ref{eq:pinned}).

\paragraph{Omori children} are triggered by previous events of either class. A parent
$p$ of (released) magnitude $M_p$ is \emph{eligible} while it can still produce a child
above the catalog floor, $M_p \ge \Mmin + \dM = 6.2$. With the simulation on a regular
time grid of step $\delta t = 1$~yr, an eligible parent contributes, at integer lag
$k \in \{1,\dots,k_{\max}\}$ ($k_{\max}\delta t = T_\Omega = 30$~yr), the step-averaged
Omori--Utsu rate
\begin{equation}
  r_p(k) \;=\; \frac{K(M_p)}{\delta t}\Bigl[F(k\,\delta t) - F\bigl((k-1)\,\delta t\bigr)\Bigr],
  \qquad
  F(t) = \int_0^{t}\frac{d\tau}{(\tau + c_\Omega)^{q}},
\label{eq:omori}
\end{equation}
with exponent $q = 1$, offset $c_\Omega = 1$~day, and productivity
\begin{equation}
  K(M) \;=\; \kappa_K\, K_{\mathrm{ref}}\, 10^{\,\alpha (M - M_{\mathrm{ref}})},
  \qquad K_{\mathrm{ref}} = 0.043,\ \ \alpha = 0.8,\ \ M_{\mathrm{ref}} = 6,
\label{eq:K}
\end{equation}
and the productivity multiplier $\kappa_K = 10$ that defines this configuration.
Integrating the kernel over the step (rather than point-sampling it) preserves the full
first-year productivity on the coarse grid; the expected number of direct children of a
parent is $n(M_p) = K(M_p)\,F(T_\Omega)$ with $F(T_\Omega) = \ln(1 + T_\Omega/c_\Omega)
\approx 9.30$ in years. A child's magnitude follows the \emph{full} catalog G-R law,
capped below its parent by the B\aa th gap:
\begin{equation}
  p_{\mathrm{O}}(M \mid M_p) \;\propto\; 10^{-bM},
  \qquad M \in \bigl[\Mmin,\ \min(\Mmax,\, M_p - \dM)\bigr],
  \qquad \dM = 1.2 .
\label{eq:childlaw}
\end{equation}
Children of eligible parents can themselves trigger (grandchildren), but because a child
requires $M \ge 6.2$ to be a parent and children are capped at $M_p - 1.2 \le 6.3$, the
cascade is effectively one generation deep here ($\nO = 2.5\times10^{-3}$,
Section~\ref{sec:budget}).

The total rate is $\lambda(t) = \lamL(t) + \sum_{p} r_p(t)$, the number of events per
step is Poisson, $N_{\mathrm{ev}} \sim \mathrm{Poisson}(\lambda\,\delta t)$, and each
event's origin is a categorical draw: Omori with probability
$\sum_p r_p/\lambda$ (parent chosen $\propto r_p$), loading otherwise.

\subsection{Hypocenters, rupture, and slip}
\label{sec:space}

The hypocenter of a loading event of magnitude $M$ follows the MaxEnt law
\begin{equation}
  P(i \mid M) \;\propto\; h_i\, m_i^{\gamma(M)},
  \qquad
  \gamma(M) = \gamma_{\max} - (\gamma_{\max}-\gamma_{\min})
  e^{-\alpha_s (M - \Mmin)},
\label{eq:spatial}
\end{equation}
with $\gamma_{\min} = 0$, $\gamma_{\max} = 1.5$, $\alpha_s = 0.35$: large events demand
a pool of stored moment while small events nucleate almost uniformly; the factor $h_i$
keeps nucleation off freshly shadowed patches. An Omori child nucleates on its parent's
footprint: $P(i \mid M, p) \propto h_i\, \Phi_p(i)\, m_i^{\gamma(M)}$, where
$\Phi_p \in [0,1]$ is the parent's spatial activation kernel (anisotropic exponential
decay from the ruptured elements, correlation lengths $10$~km, shared with the afterslip
module).

Given a hypocenter and magnitude, the rupture area follows the strike-slip
magnitude--area scaling $A(M) = 10^{\,M - 3.99}\ \mathrm{km^2}$ \cite{allen2017};
elements join in order of distance from the hypocenter until $A(M)$ is reached, receive
an exponentially tapered heterogeneous slip pattern scaled to the target moment
$10^{1.5M+9.05}$~N\,m, and the event is \emph{clipped to the available deficit}: slip is
scaled down uniformly if the target exceeds the deficit stored on the ruptured elements
and clamped per element, and the released magnitude is recomputed from the slip actually
applied (both the nominal draw and the released magnitude are recorded; at
$\Dref = 250$~yr clipping affects $0.4\%$ of events). Events with $M \ge 7$ additionally
start an afterslip sequence that releases deficit aseismically on the rupture halo
(budget capped at the coseismic moment); the aseismic release fraction is
$\fas = 0.2$ of loading.

\section{Closing the moment budget over the cascade}
\label{sec:budget}

The new calibration is the piece that makes the split self-consistent. Let
$p_{\mathrm{L}}$ be the loading law (\ref{eq:loadlaw}) and define, on the child law
(\ref{eq:childlaw}),
\begin{align}
  \nL &= \mathbb{E}_{p_{\mathrm{L}}}\!\left[\,n(M)\,\mathbf{1}\{M \ge \Mmin+\dM\}\right]
        && \text{children per loading event,} \nonumber\\[2pt]
  \EO &= \frac{\mathbb{E}_{p_{\mathrm{L}}}\!\left[n(M)\,\bar m_{\mathrm{child}}(M)\right]}
              {\mathbb{E}_{p_{\mathrm{L}}}\!\left[n(M)\right]},
        \qquad
        \bar m_{\mathrm{child}}(M) = \mathbb{E}\bigl[m \mid M' \sim p_{\mathrm{O}}(\cdot\mid M)\bigr],
        && \text{mean child moment,} \nonumber\\[2pt]
  \nO &= \text{expected children of a child (same construction, one level down).}
        &&
\label{eq:defs}
\end{align}
Writing $\bar\lambda_{\mathrm{L}}$ and $\bar\lambda_{\mathrm{O}} =
\bar\lambda_{\mathrm{L}}\,\nL/(1-\nO)$ for the long-run loading and child rates, the
seismic moment balance $(1-\fas)\Ltot = \bar\lambda_{\mathrm{L}}\EL +
\bar\lambda_{\mathrm{O}}\EO$ gives the loading rate in closed form:
\begin{equation}
  \boxed{\;
  \bar\lambda_{\mathrm{L}}
  \;=\;
  \frac{(1-\fas)\,\Ltot}{\EL + \dfrac{\nL\,\EO}{1-\nO}}
  \;}
\label{eq:budget}
\end{equation}
--- the moment released per loading event \emph{including its whole sequence} replaces
$\bar m$ in Eq.~(\ref{eq:pinned}). For this configuration ($\kappa_K = 10$):
\[
  \EL = 3.571\times10^{8}\ \mathrm{m^3},\quad
  \nL = 8.49,\quad
  \EO = 3.54\times10^{6}\ \mathrm{m^3} = 0.010\,\EL,\quad
  \nO = 2.5\times10^{-3},
\]
\[
  \bar\lambda_{\mathrm{L}} = 0.158\ \mathrm{yr^{-1}},\qquad
  \bar\lambda_{\mathrm{O}} = 1.35\ \mathrm{yr^{-1}},\qquad
  \bar\lambda_{\mathrm{tot}} = \bar\lambda_{\mathrm{L}}\Bigl(1 + \tfrac{\nL}{1-\nO}\Bigr)
  = 1.51\ \mathrm{yr^{-1}} .
\]
The sequences carry $\bar\lambda_{\mathrm{O}}\EO/[(1-\fas)\Ltot] = 7.8\%$ of the seismic
moment flux but $\nL/(1+\nL) \approx 90\%$ of the events. Equation~(\ref{eq:budget}) is
evaluated once at start-up by quadrature on the two magnitude grids; the earlier
approximation that treated aftershocks as releasing the same mean moment as background
events ($\bar\lambda_{\mathrm{L}} \propto 1-n_b$) under-drives the fault by
$\sim27\%$ and lets the reservoir drift $\sim30\%$ high, so the cascade form is required
whenever the split is on. With $\bar\lambda_{\mathrm{L}}$ fixed,
$\lambda_0 = \bar\lambda_{\mathrm{L}}/\langle h\rangle_{\mathrm{ref}} = 0.264\
\mathrm{yr^{-1}}$, and the controller supplies only a trim.

\section{The sequence-generation algorithm}
\label{sec:algo}

\begin{algorithm}[t]
\caption{One realization of the moment-budget/Omori-split model
(time step $\delta t$, duration $T$; all random draws from one seeded generator).}
\label{alg:kes}
\begin{algorithmic}[1]
\State \textbf{initialize} $s_i \gets v_i T_{\mathrm{ref}}$;\ \ $\xi_i \gets$ steady-state
       shadow population;\ \ compute $\EL,\nL,\EO,\nO$ and
       $\bar\lambda_{\mathrm{L}}$ by Eq.~(\ref{eq:budget});\ \
       $\lambda_0 \gets \bar\lambda_{\mathrm{L}} / \langle h\rangle_{\mathrm{ref}}$;\ \ $c \gets 1$
\For{each step $t \gets \delta t, 2\delta t, \dots, T$}
  \State advance active afterslip sequences (aseismic deficit release, clipped locally)
  \State $s_i \mathrel{+}= v_i\,\delta t$;\ \ relax shadow
         $\xi_i \gets 1 + (\xi_i - 1)e^{-\delta t / t_a}$;\ \ $h_i \gets 1/\xi_i$
  \State \textbf{parent scan:} for each past event with lag $k \le k_{\max}$ and
         $M_p \ge \Mmin + \dM$, collect $r_p(k)$ by Eq.~(\ref{eq:omori});
         free the $\Phi$ kernels of expired parents
  \State $\lamL \gets \lambda_0\, c\, \langle h\rangle\, (\Dres/\Dref)^{\nu}$;\ \
         $\lambda \gets \lamL + \textstyle\sum_p r_p$;\ \
         $N_{\mathrm{ev}} \sim \mathrm{Poisson}(\lambda\,\delta t)$
  \For{each of the $N_{\mathrm{ev}}$ events}
    \If{$U(0,1) < \sum_p r_p / \lambda$} \Comment{\emph{Omori child}}
       \State parent $p \propto r_p$;\ \ $M \sim$ G-R on
              $[\Mmin,\ \min(\Mmax, M_p - \dM)]$;\ \ weights $w_i \gets h_i \Phi_p(i)$
    \Else \Comment{\emph{loading event}}
       \State $M \sim$ G-R on $[\ML, \Mmax]$;\ \ weights $w_i \gets h_i$
    \EndIf
    \State hypocenter $i^\ast \sim P(i) \propto w_i\, m_i^{\gamma(M)}$
    \State grow rupture to area $10^{M - 3.99}\ \mathrm{km^2}$; taper slip; scale to
           $10^{1.5M + 9.05}$ N\,m; \textbf{clip} to available deficit; recompute
           released $M$
    \State $s_i \mathrel{-}= u_i$;\ \ shadow step
           $\ln\xi_i \mathrel{+}= f\,u_i/(v_i t_a)$ on ruptured elements
    \State record event (origin, parent, nominal and released $M$, hypocenter, slip);
           if $M \ge 7$ start afterslip and store $\Phi$; else store $\Phi$ if $M \ge \Mmin$
  \EndFor
  \State controller: $c \gets \mathrm{clip}\bigl(c + g\,(\Dres/\Dref - 1)\,\delta t,\ 0.5,\ 2\bigr)$
\EndFor
\end{algorithmic}
\end{algorithm}

Algorithm~\ref{alg:kes} lists one realization. Three implementation details matter for
reproducibility and stationarity. (i)~\emph{Within-step ordering}: multiple events in
one step act on a working copy of the deficit field that is decremented event by event,
so a large event immediately shadows and depletes the field seen by the next draw.
(ii)~\emph{Step-integrated triggering}: Eq.~(\ref{eq:omori}) delivers the whole
first-year productivity at lag one; point-sampling the kernel on a 1-yr grid would
silently discard $\sim57\%$ of it. (iii)~\emph{Kernel lifetime}: a parent's spatial
kernel $\Phi_p$ is kept in memory exactly while the parent is inside the Omori window
and freed afterwards. Every random quantity (Poisson counts, origin, parent, magnitude,
hypocenter, slip heterogeneity) comes from a single seeded stream, and with the split
disabled the implementation reproduces the classical model bit for bit.

\section{Emergent statistics}
\label{sec:results}

\begin{table}[t]
\centering
\small
\begin{tabular}{lcc}
\toprule
 & classical split control & this model \\
 & (loading G-R on $[5,7.5]$, $\kappa_K{=}1$) & (loading on $[6,7.5]$, $\kappa_K{=}10$) \\
\midrule
$M\ge5$ rate (yr$^{-1}$) & 1.66 & 1.55 \\
Omori-child fraction & 0.08 & 0.89 \\
loading-event rate (yr$^{-1}$) & 1.53 & 0.166 \\
geometric coupling & 0.985 & 0.985 \\
reservoir (yr of loading) & $250 \pm 58$ & $275 \pm 65$ \\
controller factor range & $[0.91, 1.16]$ & $[0.93, 1.17]$ \\
interevent CV ($M\ge5$) & 1.4 & 2.4 \\
Fano factor, 1-yr windows & 1.3 & 8.6 \\
apparent $b$ (Aki, $M\ge5$) & 1.05 & 1.32 \\
B\aa th gap (median $M_p - \max M_{\mathrm{child}}$) & 1.33 & 1.22 \\
$M\ge7$ recurrence mean / CV & 91 yr / 0.84 & 87 yr / 0.87 \\
$M\ge5$ counts, 10 yr before / after largest event & 16.2 / 11.9 & 15.1 / 54.5 \\
$M\ge5$ rate ratio $[-10,-1]$ / $[+1,+10]$ & 1.06 / 1.30 & 0.85 / 2.86 \\
$M\ge6.5$ local rate ratio $[+1,+50]$ & 0.40 & 0.23 \\
\bottomrule
\end{tabular}
\caption{Ensemble statistics, five runs $\times$ 3000 yr each (seeds 42--46).
``Local'' restricts to $\pm25$~km of the mainshock. Windows are integer 1-yr grid lags
around $M\ge7$ mainshocks except the ``largest event'' row (largest event of each
1000-yr block).}
\label{tab:stats}
\end{table}

\begin{figure}[t]
  \centering
  \includegraphics[width=\textwidth]{lollipop_load6_comparison.png}
  \caption{The same 1000-yr window construction for (a) the split control with the
  classical full-range loading law, (b) the loading floor $\ML = 6$ at classical
  productivity ($\kappa_K = 0.9$; a sparse, depinning-like catalog at
  $0.29\ M{\ge}5\,\mathrm{yr^{-1}}$), and (c) this model ($\kappa_K = 10$): the classical
  total rate, reorganized into sequences.}
  \label{fig:compare}
\end{figure}

Table~\ref{tab:stats} compares five-seed ensembles of this model against the split
control that keeps the classical full-range loading law. The conservation diagnostics
are indistinguishable --- coupling $0.985$, reservoir within $10\%$ of $\Dref$,
controller trim within $\pm17\%$ --- while the catalog texture changes completely: the
interevent coefficient of variation rises from $1.4$ to $2.4$ and the one-year Fano
factor from $1.3$ to $8.6$ (strong clustering), the decade after a large event carries
$2.9\times$ the long-run $M\ge5$ rate while the decade before carries $0.85\times$, and
the counts around the largest event of each millennium go from $16/12$
(pre/post, classical) to $15/55$. Two magnitude-composition signatures are worth
flagging. First, the apparent $b$-value of the \emph{mixed} catalog is not the
prescribed $b=1$: the population is bimodal (loading events at $M\ge6$, children at
$M \le 6.3$), and an Aki fit across the mixture returns $b \approx 1.3$; the prescribed
law is recovered within each class separately. Second, small-event quiescence deepens
with proximity: within $25$~km of a mainshock the $M\ge6.5$ rate in the following half
century is $0.23\times$ its long-run value --- the stress-shadow memory acting on top of
the sequence organization. B\aa th's law holds by construction (median gap $1.22$; no
child ever exceeds its parent).

What this configuration does \emph{not} do: the loading events themselves remain
memoryless in size ($\rho(\Dres, M \mid M\ge7) \approx 0$), because the deficit-tilted
magnitude law of the companion study is switched off here ($\mu = 0$); the two
mechanisms compose freely if reservoir-timed large events are also wanted. The post/pre
asymmetry of $\sim3.6$ at ten years is set by the productivity multiplier and the Omori
exponent, not by the split itself; pushing toward the $\sim17{:}1$ asymmetry of the
depinning comparison model is a matter of raising $\alpha$ (concentrating productivity
under the largest parents) and belongs to a separate calibration exercise.

\section{Reproducibility}
\label{sec:repro}

The model is implemented in the KES repository (\texttt{magnitude\_law.py},
\texttt{temporal\_prob.py}, \texttt{simulator.py}); with all new keys at their defaults
the code reproduces the classical model bit for bit. This configuration is exactly:
\begin{center}
\small
\begin{tabular}{ll}
\toprule
\texttt{omori\_split\_enabled} & \texttt{true} \\
\texttt{magnitude\_load\_M\_min} ($\ML$) & 6.0 \\
\texttt{omori\_bath\_K\_scale} ($\kappa_K$) & 10.0 \\
\texttt{omori\_bath\_dM} ($\dM$) & 1.2 \\
(all other parameters) & repository defaults \\
\bottomrule
\end{tabular}
\end{center}
The ensemble and every figure and number in this note regenerate with:
\begin{quote}\small\ttfamily
python run\_ensemble.py --vary\_seed --n\_runs 5 --duration 3000 \textbackslash\\
\hspace*{1em}--set snapshot\_interval\_years=10 --set omori\_split\_enabled=true \textbackslash\\
\hspace*{1em}--set magnitude\_load\_M\_min=6 --set omori\_bath\_K\_scale=10 \textbackslash\\
\hspace*{1em}--output\_dir results/ensemble\_load6\_k10\\[4pt]
python plot\_foreshocks.py --input\_dir results/ensemble\_load6\_k10 \textbackslash\\
\hspace*{1em}--m\_mainshock 7.0 --m\_targets 5,6.5 --window 50\\
python plot\_lollipop.py --input\_dir results/ensemble\_load6\_k10\\
python plot\_gr\_omori.py --input\_dir results/ensemble\_load6\_k10
\end{quote}
A five-seed, 3000-yr ensemble runs in about three minutes on a laptop.

\begin{thebibliography}{9}\small
\bibitem{allen2017} Allen, T.~I., and G.~P. Hayes (2017), Alternative rupture-scaling
relationships for subduction interface and other offshore environments,
\emph{Bull. Seismol. Soc. Am.}, 107, 1240--1253.
\bibitem{bath1965} B\aa th, M. (1965), Lateral inhomogeneities of the upper mantle,
\emph{Tectonophysics}, 2, 483--514.
\bibitem{dieterich1994} Dieterich, J. (1994), A constitutive law for rate of earthquake
production and its application to earthquake clustering,
\emph{J. Geophys. Res.}, 99, 2601--2618.
\bibitem{kagan2002} Kagan, Y.~Y. (2002), Seismic moment distribution revisited: I.
Statistical results, \emph{Geophys. J. Int.}, 148, 520--541.
\bibitem{reasenberg1989} Reasenberg, P.~A., and L.~M. Jones (1989), Earthquake hazard
after a mainshock in California, \emph{Science}, 243, 1173--1176.
\bibitem{utsu1961} Utsu, T. (1961), A statistical study on the occurrence of
aftershocks, \emph{Geophys. Mag.}, 30, 521--605.
\end{thebibliography}

\end{document}
